% ======================================================================
\section{Implementation and Experimental Setup}
\label{sec:exp}
% ======================================================================
\subsection{Datasets and Protocol}
\textbf{UNSW-NB15}~\cite{moustafa2015unsw}: we use the official partition
(175{,}341 training / 82{,}332 test flows; 9 attack families; 39 numeric
and 3 categorical attributes, i.e.\ 42 tokens). 15\% of the training part
is held out for validation, and the test set is used only for final
reporting. Its attack share differs from training (55.1\% vs.\ 68.1\%),
which makes this the harder, more realistic setting.
\textbf{CICIDS2017}~\cite{sharafaldin2018cicids}: we use the publicly
released stratified subset of Yang et al.~\cite{yang2022mthids}
(56{,}661 flows, 77 features, 6 grouped attack classes: DoS, PortScan,
brute force, web attack, bot and infiltration; 60 informative numeric
features). We split it 70/15/15 (39{,}662/8{,}499/8{,}500 flows) with
stratification. Given known labelling issues~\cite{engelen2021troubleshooting},
we treat it as a second, easier benchmark.
Duplicate and constant columns (decided on training data only) are removed,
infinite values are imputed with training medians, and categorical
attributes keep their 31 most frequent levels plus an ``other'' bucket.
SMOTE-NC~\cite{chawla2002smote} raises every attack family to at least 2\% of
the \emph{training} split only. Every configuration is run with three seeds
(0, 1, 2); the seed changes the train/validation split, SMOTE, weight
initialization and the DP noise. We report mean${}\pm{}$std.


\subsection{Baselines}
All baselines receive the same privatized inputs:
(i) Isolation Forest~\cite{liu2008iforest}, trained on benign flows only, with
the threshold chosen on validation F1;
(ii) Random Forest~\cite{breiman2001rf} (300 trees);
(iii) XGBoost~\cite{chen2016xgboost}, configured exactly like the HALX head;
(iv) a two-layer MLP (256 units);
(v) CNN-LSTM in the style of~\cite{halbouni2022cnnlstm} (two Conv1d layers
followed by an LSTM over feature tokens);
(vi) the HALX BiLSTM without attention and without the XGBoost head; and
(vii) a feature-token Transformer (FT-Transformer style~\cite{gorishniy2021revisiting},
two layers, 4 heads, $d=64$), which stands in for the transformer IDS
family~\cite{wu2022rtids,manocchio2024flowtransformer}. The deep baselines
use the same optimizer, schedule and early-stopping rule as HALX.

\subsection{Implementation}
HALX uses PyTorch~2, XGBoost~3 and scikit-learn; the privacy layer
(Mondrian tree, zCDP standardizer) is ${\approx}150$ lines of NumPy. All
runs used a 4-core x86 CPU (16\,GB RAM, no GPU), at most 15 epochs (8 on
UNSW-NB15) and patience 3. We report accuracy, precision, recall, F1 and
FPR for the binary decision at $\tau=0.5$, ROC-AUC, and macro-F1 (M-F1)
over attack families.
